% Figure 3: reading over a network. Every value is read from generated/fignet_*.tsv and generated/figr2.tsv (written by
% `python -m s1bench report`, network_section, netsweep_section and r2_section). More panels: fig_supp_net.tex.
% Panels on a 2 x 2 grid: .s1 in teal, every other format in gray; the lines of (b) are labeled at their right end.
\begin{figure*}[tbp]
\centering
\pgfmathsetlengthmacro{\netW}{0.33\linewidth}   % axis width
\pgfmathsetlengthmacro{\netH}{0.24\linewidth}   % axis height
\pgfmathsetlengthmacro{\netG}{1.6cm}            % horizontal gap (y labels)
\pgfmathsetlengthmacro{\netV}{1.5cm}            % vertical gap (x labels and titles)
\begin{tikzpicture}
\pgfplotsset{netax/.style={scale only axis, width=\netW, height=\netH, axis x line*=bottom, axis y line*=left}}
% ---------------- (a) measured training loaders against link bandwidth ----------------
\begin{axis}[netax, name=nb, at={(0,0)}, anchor=north west,
  title={a\quad Measured training loaders},
  xmode=log, ymode=log, xmin=0.03, xmax=30, ymin=50, ymax=400000,
  xlabel={link bandwidth (Gbit/s)}, ylabel={cells per second into the model}, xtick={0.1,1,10}, xticklabels={0.1,1,10}]
\addplot[c-s1, very thick, mark=*, mark size=1.8pt, discard if not={loader}{s1_stream}] table[x=gbit, y=cells_per_s] {\gen{fignet_loaders.tsv}}
  node[pos=0.75, anchor=north, font=\sffamily\tiny\bfseries, text=c-s1!85!black, inner sep=3pt] {\sone};
\addplot[c-zarr, thick, mark=square*, mark size=1.4pt, discard if not={loader}{anndata_zarr_lazy}] table[x=gbit, y=cells_per_s] {\gen{fignet_loaders.tsv}}
  node[pos=0.75, anchor=north, font=\sffamily\tiny, text=c-zarr!70!black, inner sep=3pt] {AnnData Zarr};
\addplot[c-soma, thick, mark=triangle*, mark size=1.6pt, discard if not={loader}{soma_ml}] table[x=gbit, y=cells_per_s] {\gen{fignet_loaders.tsv}}
  node[pos=0.75, anchor=north, font=\sffamily\tiny, text=c-soma!70!black, inner sep=3pt] {TileDB-SOMA ML};
\addplot[c-s1, dashed, no marks] table[x=gbit, y=net] {\gen{fignet_ceilings.tsv}}
  node[pos=0.2, sloped, above, font=\sffamily\tiny, text=c-s1] {link capacity, in \sone{} cells};
\end{axis}
% ---------------- (b) delivered rate against link bandwidth ----------------
\begin{axis}[netax, name=na, at={($(nb.north east)+(\netG,0)$)}, anchor=north west,
  title={b\quad What a link delivers, by format},
  xmode=log, ymode=log, xmin=0.01, xmax=100, ymin=0.3, ymax=1500,
  xlabel={link bandwidth (Gbit/s)}, ylabel={million non-zeros per second},
  xtick={0.01,0.1,1,10,100}, xticklabels={0.01,0.1,1,10,100}, every axis plot/.append style={no marks}, clip mode=individual,
  after end axis/.code={\input{\gen{fignet_labels.tex}}}]   % each line labeled at the fastest link
\addplot[c-s1, line width=1.6pt] table[x=gbit, y=s1] {\gen{fignet_curves.tsv}};
\addplot[c-zarr, thick] table[x=gbit, y=zarr] {\gen{fignet_curves.tsv}};
\addplot[c-zarr3, thick] table[x=gbit, y=zarr3] {\gen{fignet_curves.tsv}};
\addplot[c-soma, thick] table[x=gbit, y=soma] {\gen{fignet_curves.tsv}};
\addplot[c-h5ad_gzip, thick] table[x=gbit, y=h5ad_gzip] {\gen{fignet_curves.tsv}};
\addplot[c-csc_i32, thick] table[x=gbit, y=csc_i32] {\gen{fignet_curves.tsv}};
\addplot[c-csc_zstd19, thick] table[x=gbit, y=csc_zstd19] {\gen{fignet_curves.tsv}};
\addplot[c-tenx_h5, thick] table[x=gbit, y=tenx_h5] {\gen{fignet_curves.tsv}};
\addplot[c-bpcells, thick] table[x=gbit, y=bpcells] {\gen{fignet_curves.tsv}};
\addplot[c-parquet_coo, thick] table[x=gbit, y=parquet_coo] {\gen{fignet_curves.tsv}};
\addplot[c-rds_xz, thick] table[x=gbit, y=rds_xz] {\gen{fignet_curves.tsv}};
\addplot[c-rds_gz, thick] table[x=gbit, y=rds_gz] {\gen{fignet_curves.tsv}};
\addplot[c-mtx_gz, thick] table[x=gbit, y=mtx_gz] {\gen{fignet_curves.tsv}};
\addplot[c-loom, thick] table[x=gbit, y=loom] {\gen{fignet_curves.tsv}};
\addplot[c-h5ad_lzf, thick] table[x=gbit, y=h5ad_lzf] {\gen{fignet_curves.tsv}};
\addplot[c-h5ad_none, thick] table[x=gbit, y=h5ad_none] {\gen{fignet_curves.tsv}};
\addplot[c-s1, line width=1.6pt, forget plot] table[x=gbit, y=s1] {\gen{fignet_curves.tsv}};   % on top
\addplot[c-s1, densely dotted, forget plot] table[x=gbit, y=s1_net] {\gen{fignet_curves.tsv}}
  node[pos=0.22, sloped, above, font=\sffamily\tiny, text=c-s1] {limited by the link};
\addplot[c-s1, densely dotted, forget plot] table[x=gbit, y=s1_dec] {\gen{fignet_curves.tsv}}
  node[pos=0.02, anchor=south west, font=\sffamily\tiny, text=c-s1] {limited by the decoder};
\end{axis}
% ---------------- (c) sustained rate against added round-trip time ----------------
\begin{axis}[netax, name=nc, at={($(nb.south west)+(0,-\netV)$)}, anchor=north west,
  title={c\quad Rate against round-trip time, 1\,Gbit/s},
  xmode=log, ymode=log, xmin=0.4, xmax=260, ymin=100, ymax=2e5,
  xtick={0.5,5,10,30,100,200}, xticklabels={0,5,10,30,100,200}, ytick={1e2,1e3,1e4,1e5},
  xlabel={added round-trip time (ms)}, ylabel={cells per second into the model}, clip=false]
\addplot[c-s1, very thick, mark=*, mark size=1.6pt] table[x=rtt, y=s1_stream_cps] {\gen{fignet_rtt.tsv}}
  node[pos=0, anchor=south west, font=\sffamily\tiny\bfseries, text=c-s1!85!black, inner sep=3pt] {\sone};
\addplot[c-zarr, thick, mark=square*, mark size=1.3pt] table[x=rtt, y=anndata_zarr_lazy_cps] {\gen{fignet_rtt.tsv}}
  node[pos=0, anchor=south west, font=\sffamily\tiny, text=c-zarr!70!black, inner sep=3pt] {AnnData Zarr};
\addplot[c-zarr3, thick, mark=diamond*, mark size=1.8pt] table[x=rtt, y=zarr3_shard_cps] {\gen{fignet_rtt.tsv}}
  node[pos=0, anchor=south west, font=\sffamily\tiny, text=c-zarr3!70!black, inner sep=3pt] {Zarr v3, sharded};
\addplot[c-soma, thick, mark=triangle*, mark size=1.6pt] table[x=rtt, y=soma_ml_cps] {\gen{fignet_rtt.tsv}}
  node[pos=0, anchor=south west, font=\sffamily\tiny, text=c-soma!70!black, inner sep=3pt] {TileDB-SOMA ML};
\end{axis}
% ---------------- (d) a real object store ----------------
\begin{axis}[netax, name=nd, at={($(na.south west)+(0,-\netV)$)}, anchor=north west,
  title={d\quad A real cloud object store},
  xmode=log, ymode=log, xmin=0.3, xmax=300, ymin=150, ymax=6e5,
  xlabel={seconds to the first minibatch}, ylabel={cells per second, sustained}, xtick={1,10,100}, xticklabels={1,10,100},
  lbl/.style={nodes near coords, point meta=explicit symbolic, visualization depends on={value \thisrow{anchor}\as\lblanchor},
    every node near coord/.append style={font=\sffamily\tiny, anchor=\lblanchor, text=ink2, inner sep=2pt}}]
\fill[c-s1!7] (axis cs:0.3,1e4) rectangle (axis cs:1,6e5);
\node[font=\sffamily\tiny, text=c-s1!80!black, anchor=north west, align=left] at (axis cs:0.31,5.5e5) {fast to start\\and to stream};
% lines of equal time to one epoch through the .s1 point, the downloaded h5ad file and AnnData Zarr; labels from
% generated/figr2_epoch_labels.tex, horizontal, below each line's flat part from 2 s
\addplot[c-s1, densely dashed, thick, no marks, discard if not={curve}{s1}] table[x=first, y=cps] {\gen{figr2_epoch.tsv}};
\addplot[c-h5ad_gzip, densely dashed, thick, no marks, discard if not={curve}{h5ad}] table[x=first, y=cps] {\gen{figr2_epoch.tsv}};
\addplot[c-zarr, densely dashed, thick, no marks, discard if not={curve}{zarr}] table[x=first, y=cps] {\gen{figr2_epoch.tsv}};
\input{\gen{figr2_epoch_labels.tex}}
\pgfplotsinvokeforeach{h5ad_gzip,zarr,zarr3,s1}{
  \addplot[only marks, lbl, mark=*, mark size=2.2pt, c-#1, discard if not={fmt}{#1}] table[x=first, y=cps, meta=label] {\gen{figr2.tsv}};
}
\end{axis}
\end{tikzpicture}
\caption{\textbf{\sone{} delivers the most cells over every link up to 1\,Gbit/s on four cores, including from a real
cloud object store.} Every loader decodes what it fetches into PyTorch sparse tensors; \sone{} in teal, every other format in gray.
\textbf{(a)} Cells per second delivered into a PyTorch model from \nm{fig1b-dataset} through an emulated object store at
three link speeds (1\,ms latency); dashed, the link's capacity in \sone{} cells.
\textbf{(b)} Non-zeros a link of a given bandwidth delivers in each format, from its median size and decode rate on
\nm{threadsHeadline} cores (Eq.~\ref{eq:reff}). Each curve rises with the link until the decoder becomes the limit;
\sone{} carries the most non-zeros over every link up to about \nm{wire-crossover-gbit}\,Gbit/s, where uncompressed arrays
catch up.
\textbf{(c)} The same store at 1\,Gbit/s with a round-trip delay added to every request, \nm{ns-dataset}. The round trip
is the time from sending a request to receiving the first byte of its answer, set by distance and the storage service, not
by computing: a few milliseconds within a data center, tens of milliseconds from cloud storage over the internet. A loader
that waits for one request before sending the next pays it on every request; \sone{} with the \texttt{singlet}
\nm{rel-singlet_py} defaults keeps 32 range requests in flight and requests the next four minibatches ahead.
\textbf{(d)} Cloudflare R2 read over the public internet from a cluster node (about \nm{r2-hlca_core-rtt}\,ms per
request), \nm{fig1b-dataset}, median of three seeds; the h5ad file is downloaded whole before the first minibatch. Dashed:
loaders on a line deliver one epoch (\nm{r2-epoch-cells} cells) in the same time, through the \sone{}, downloaded h5ad and
AnnData Zarr points; the downloaded file catches up with \sone{} only after about \nm{r2-epoch-crossover} epochs, if it fits
in memory. Costs per
epoch: Supplementary Table~\ref{S-tab:r2}; latency distributions, round trips and block size: Supplementary
Fig.~\ref{S-fig:snet}.}
\label{fig:net}
\end{figure*}
